\documentclass[11pt]{article}

\usepackage[margin=1in]{geometry}
\usepackage{fontspec}
\usepackage{polyglossia}
\setmainlanguage{russian}
\setotherlanguage{english}
\setmainfont{Times New Roman}
\newfontfamily\cyrillicfont{Times New Roman}[Script=Cyrillic]
\newfontfamily\cyrillicfonttt{Courier New}[Script=Cyrillic]
\usepackage{amsmath}
\usepackage{amssymb}
\usepackage{booktabs}
\usepackage{array}
\usepackage{tabularx}
\usepackage{xurl}
\usepackage{hyperref}
\usepackage{authblk}
\newcolumntype{P}[1]{>{\raggedright\arraybackslash}p{#1}}
\newcolumntype{Y}{>{\raggedright\arraybackslash}X}

\title{Пятиуровневый аудит конечнополевой параметризации e-5-137}

\author{И. Михайлов}
\affil{Независимый исследователь, Москва, Россия}

\date{30 мая 2026}

\begin{document}

\maketitle

\begin{abstract}
В работе фиксируется конечнополевая параметризация, построенная на $e$, $\pi$,
$\alpha^{-1}$, $N=5$, $D=26$ и арифметике GF(137).  Сначала дан атлас
стандартных формул механики, электродинамики, термодинамики, теории
относительности, квантовой механики, ядерной физики и космологии.  Затем
описан пятиуровневый аудит: MeV-масштаб ядерных дефектов масс, eV-масштаб
лептонных соотношений, молекулярный и семантический уровень восстановления,
макроскопический предел идеального газа и космологические проверки.  Модель
является феноменологической параметризацией.  Она не заменяет Стандартную
модель, КХД, статистическую механику или $\Lambda$CDM.  Каждый раздел поэтому
содержит либо численный контрольный результат, либо отрицательный результат.
\end{abstract}

\section{Область применения}

В проекте используются постоянные
\begin{equation}
    N=5,\qquad D=26,\qquad p=137.
\end{equation}
Арифметика конечного поля задается вычетами
\begin{equation}
    \mathbb{F}_{137}\simeq\mathbb{Z}/137\mathbb{Z}.
\end{equation}
Основной безразмерный оператор сжатия:
\begin{equation}
    q=\frac{N(N-2)}{e^4\pi^3}
      =0.008860611874290873.
\end{equation}
Фазовый сдвиг:
\begin{equation}
    \delta\phi=\cos(1/N)-\cos(2/N)
              =0.05900558383835652.
\end{equation}
Порог, который используется в целочисленных ядрах и каналах восстановления:
\begin{equation}
    I_5=2(D-N)=42.
\end{equation}

Задача текста не состоит в том, чтобы выводить физику из совпадений чисел.
Задача состоит в том, чтобы явно указать формулы, данные, остатки и области
отказа.  Формула остается в тексте только тогда, когда для нее указан масштаб
применимости и численный тест.

Эта работа входит в связанную серию репозиториев-аудитов e-5-137 и GF(137).
Репозитории намеренно используют общие обозначения, приемы реализации и
правила воспроизводимости.  Область данного архива ограничена аудитом
конечного поля, активной матрицей из пяти тестов и двумя версиями препринта
для этой статьи.  Результаты соседних репозиториев следует цитировать по их
собственным архивным релизам или DOI.

\section{Общие операторы и структуры данных}

В вычислительной части используются три простые структуры.

Первая структура -- байтовый вектор вычетов:
\begin{equation}
    x,w,b\in\mathbb{F}_{137}.
\end{equation}
Целочисленное ядро инференса считает
\begin{equation}
    R=(XW+b)\bmod 137,\qquad A_{ij}=\mathbf{1}_{R_{ij}\ge42}.
\end{equation}
Это инженерный примитив для компактного инференса, а не новая теория нейронных
сетей.

Вторая структура -- канал из $26$ осей.  Символ $s$ реплицируется как
\begin{equation}
    r_a=s+m_a(k)\pmod{137},\qquad a=0,\ldots,25,
\end{equation}
где $m_a(k)$ -- детерминированная маска оси.  Восстановление выбирает наиболее
частый вычет:
\begin{equation}
    \hat{s}
    =
    \operatorname*{arg\,max}_{v\in\mathbb{F}_{137}}
    \#\{a:r_a-m_a(k)\equiv v\pmod{137}\}.
\end{equation}
Если повреждено не более $10$ осей, остаются как минимум $16$ согласованных
осей:
\begin{equation}
    26-10=16.
\end{equation}
Это фиксированный бюджет восстановления для тестовой схемы.

Третья структура -- контрольный тест семантического сжатия.  Объект длиной
$120{,}000$ токенов отображается в
\begin{equation}
    V_{\rm ref}=117
\end{equation}
семантических векторов:
\begin{equation}
    \frac{120{,}000}{117}=1025.6410256410256
    \quad {\rm токенов/вектор}.
\end{equation}
Это синтетический тест из тестовой матрицы.  Он не является общей теоремой
о lossless-сжатии.

\section{Атлас формул по физическим масштабам}

Этот раздел задает справочную сетку стандартной физики.  Конечнополевая
параметризация не заменяет эти формулы.  Она используется только как слой
аудита, если явно задана дискретная модель, канал сжатия или численный
контрольный тест.

\subsection{Размерности и законы сохранения}

Первый фильтр -- размерности:
\begin{equation}
    [v]=LT^{-1},\qquad
    [a]=LT^{-2},\qquad
    [F]=MLT^{-2},\qquad
    [E]=ML^2T^{-2}.
\end{equation}
Для замкнутых систем используются стандартные сохранения:
\begin{equation}
    \frac{dE}{dt}=0,\qquad
    \frac{dp}{dt}=0,\qquad
    \frac{dL}{dt}=0.
\end{equation}
Если выражение нарушает размерности, оно не рассматривается как физическая
формула.

\subsection{Кинематика и механика Ньютона}

\begin{equation}
    v=\frac{dx}{dt},\qquad
    a=\frac{dv}{dt}=\frac{d^2x}{dt^2}.
\end{equation}
Для постоянного ускорения:
\begin{align}
    x(t)&=x_0+v_0t+\frac{1}{2}at^2,\\
    v(t)&=v_0+at,\\
    v^2&=v_0^2+2a(x-x_0).
\end{align}
Динамика:
\begin{equation}
    F=ma,\qquad p=mv,\qquad J=\int F\,dt=\Delta p.
\end{equation}
Работа и энергия:
\begin{equation}
    W=\int F\,dx,\qquad
    K=\frac{1}{2}mv^2,\qquad
    U_g=mgh,\qquad
    U_s=\frac{1}{2}kx^2.
\end{equation}
Лагранжев язык:
\begin{equation}
    L=T-V,\qquad
    \frac{d}{dt}\frac{\partial L}{\partial \dot{q}}
    -\frac{\partial L}{\partial q}=0.
\end{equation}
Слой e-5-137 здесь может использоваться только как целочисленное
представление симулятора.

\subsection{Гравитация и орбиты}

\begin{equation}
    F=G\frac{m_1m_2}{r^2},\qquad
    g=\frac{GM}{r^2},\qquad
    \Phi=-\frac{GM}{r}.
\end{equation}
Для круговой орбиты и скорости убегания:
\begin{equation}
    v_{\rm orb}=\sqrt{\frac{GM}{r}},
    \qquad
    v_{\rm esc}=\sqrt{\frac{2GM}{r}}.
\end{equation}
Третий закон Кеплера:
\begin{equation}
    T^2=\frac{4\pi^2}{GM}a^3.
\end{equation}
Релятивистское расширение не выводится из этой ньютоновской части.  Это
отдельная метрическая теория.

\subsection{Колебания, волны и оптика}

\begin{equation}
    m\ddot{x}+kx=0,\qquad
    \omega=\sqrt{\frac{k}{m}},\qquad
    T=\frac{2\pi}{\omega}.
\end{equation}
Волновое уравнение:
\begin{equation}
    \frac{\partial^2 y}{\partial t^2}
    =
    v^2\frac{\partial^2 y}{\partial x^2},
    \qquad
    v=\lambda f.
\end{equation}
Оптика:
\begin{equation}
    n=\frac{c}{v},\qquad
    n_1\sin\theta_1=n_2\sin\theta_2,
\end{equation}
\begin{equation}
    \frac{1}{f}=\frac{1}{d_o}+\frac{1}{d_i},
    \qquad
    m=-\frac{d_i}{d_o}.
\end{equation}
Дифракционные максимумы:
\begin{equation}
    d\sin\theta=m\lambda.
\end{equation}
Эти формулы задают непрерывный волновой референс для дискретного языка фазовых
осей.

\subsection{Электричество и магнетизм}

\begin{equation}
    F=k_e\frac{q_1q_2}{r^2},\qquad
    E=\frac{F}{q},\qquad
    V=\frac{W}{q}.
\end{equation}
Цепи:
\begin{equation}
    I=\frac{dq}{dt},\qquad
    U=IR,\qquad
    P=UI=I^2R=\frac{U^2}{R}.
\end{equation}
Емкость и индуктивность:
\begin{equation}
    C=\frac{Q}{U},\qquad
    E_C=\frac{1}{2}CU^2,\qquad
    E_L=\frac{1}{2}LI^2.
\end{equation}
Сила Лоренца:
\begin{equation}
    F=q(E+v\times B).
\end{equation}
Уравнения Максвелла:
\begin{align}
    \nabla\cdot E&=\frac{\rho}{\varepsilon_0},&
    \nabla\cdot B&=0,\\
    \nabla\times E&=-\frac{\partial B}{\partial t},&
    \nabla\times B&=\mu_0J+\mu_0\varepsilon_0\frac{\partial E}{\partial t}.
\end{align}
В лептонной параметризации используется $\alpha^{-1}$, но сами уравнения
Максвелла не меняются.

\subsection{Термодинамика и статистическая механика}

Идеальный газ:
\begin{equation}
    PV=nRT=Nk_BT.
\end{equation}
Первый и второй законы:
\begin{equation}
    \Delta U=Q-W,\qquad
    dS\ge \frac{\delta Q}{T}.
\end{equation}
Нагрев и фазовый переход:
\begin{equation}
    Q=mc\Delta T,\qquad
    Q=\lambda m.
\end{equation}
КПД Карно:
\begin{equation}
    \eta=1-\frac{T_c}{T_h}.
\end{equation}
Статистическая форма:
\begin{equation}
    S=k_B\ln\Omega,\qquad
    p_i=\frac{e^{-\beta E_i}}{Z},\qquad
    Z=\sum_i e^{-\beta E_i}.
\end{equation}
GF(137) на этом уровне используется только как конечный микроскопический
алфавит.  При случайном перемешивании и большом числе степеней свободы он
переходит к обычному распределению Больцмана.

\subsection{Специальная теория относительности}

\begin{equation}
    \gamma=\frac{1}{\sqrt{1-v^2/c^2}}.
\end{equation}
\begin{equation}
    \Delta t'=\gamma\Delta t,\qquad
    L'=\frac{L}{\gamma}.
\end{equation}
\begin{equation}
    E=mc^2,\qquad
    E^2=p^2c^2+m^2c^4.
\end{equation}
В космологическом аудите использовался множитель, похожий на фактор замедления
времени.  Его численный вклад оказался слишком мал и зафиксирован как
отрицательный результат.

\subsection{Квантовая механика}

\begin{equation}
    E=hf=\hbar\omega,\qquad
    p=\hbar k=\frac{h}{\lambda}.
\end{equation}
\begin{equation}
    \Delta x\,\Delta p\ge\frac{\hbar}{2}.
\end{equation}
Уравнение Шредингера:
\begin{equation}
    i\hbar\frac{\partial\psi}{\partial t}=\hat{H}\psi,
    \qquad
    \hat{H}=-\frac{\hbar^2}{2m}\nabla^2+V.
\end{equation}
Для водородоподобных уровней:
\begin{equation}
    E_n=-13.6~{\rm eV}\frac{Z^2}{n^2}.
\end{equation}
В проекте величина $5e~{\rm eV}=13.591409142295225~{\rm eV}$ используется как
численный якорь рядом с водородным масштабом.  Это не вывод КЭД.

\subsection{Ядерная физика и частицы}

\begin{equation}
    E_b=\Delta mc^2,\qquad
    R=R_0A^{1/3}.
\end{equation}
Радиоактивный распад:
\begin{equation}
    N(t)=N_0e^{-\lambda t},\qquad
    T_{1/2}=\frac{\ln2}{\lambda}.
\end{equation}
Энергия реакции:
\begin{equation}
    Q=(m_{\rm initial}-m_{\rm final})c^2.
\end{equation}
Ядерный расчет проекта проверяется против $E_b=\Delta mc^2$ и не проходит
порог $<10$ ppm для гелия, углерода и железа.

\subsection{Космология}

\begin{equation}
    v=H_0d,\qquad
    1+z=\frac{a_0}{a}.
\end{equation}
Критическая плотность:
\begin{equation}
    \rho_c=\frac{3H^2}{8\pi G}.
\end{equation}
Уравнение Фридмана:
\begin{equation}
    H^2=\frac{8\pi G}{3}\rho-\frac{kc^2}{a^2}
    +\frac{\Lambda c^2}{3}.
\end{equation}
Параметр состояния:
\begin{equation}
    w=\frac{p}{\rho c^2}.
\end{equation}
Эффективное число релятивистских степеней свободы:
\begin{equation}
    N_{\rm eff}=3.046+\Delta N_{\rm eff}.
\end{equation}
В CMB-аудите проекта:
\begin{equation}
    \Delta N_{\rm eff}
    =
    \frac{D}{137}\delta\phi(1-q)
    =
    0.011098917626279356.
\end{equation}

\subsection{Статус формул}

\begin{table}[h]
\centering
\scriptsize
\caption{Классы формул и роль конечнополевого слоя.}
\setlength{\tabcolsep}{3pt}
\begin{tabularx}{\linewidth}{@{}P{0.13\linewidth}P{0.18\linewidth}P{0.22\linewidth}P{0.20\linewidth}Y@{}}
\toprule
Область & Формула & Закон или предел & Роль GF(137) & Статус \\
\midrule
Механика & $F=ma$, $L=T-V$ & $E,p,L$ & только численное представление & референс \\
Гравитация & $F=Gm_1m_2/r^2$ & центральное поле & нет & референс \\
Волны/оптика & $v=\lambda f$ & фазовое распространение & аналогия фазовых осей & референс \\
ЭМ & Максвелл, Лоренц & заряд и поля & $\alpha^{-1}$ только в лептонном ansatz & референс \\
Термодинамика & $PV=nRT$, $S=k_B\ln\Omega$ & максимум энтропии & конечный алфавит & усреднение \\
Квантовая & $i\hbar\partial_t\psi=\hat{H}\psi$ & унитарность & $5e$ eV как якорь & параметризация \\
Ядерная & $E_b=\Delta mc^2$ & учет массы-энергии & кодировщик дефекта масс & отрицательный результат \\
Космология & Friedmann, $N_{\rm eff}$, $w$ & расширение & CMB/DESI-аудит & условный тест/нуль \\
\bottomrule
\end{tabularx}
\end{table}

\section{Пятиуровневый физический и вычислительный аудит}

\subsection{Уровень 1: MeV, ядра и КХД}

На MeV-масштабе проверяется, предсказывают ли операторы $D=26$, $42$ и $117$
дефекты масс атомных ядер.  Проверенная форма:
\begin{equation}
    \Delta M(Z,A)
    =
    Z^2E_{\rm step}(Z)
    \left[
        1+\frac{42}{137D}\left(1-\frac{13}{117}\right)
    \right].
\end{equation}
Шаг:
\begin{equation}
    E_{\rm step}(1)=5e~{\rm eV},
\end{equation}
\begin{equation}
    E_{\rm step}(Z>1)
    =
    5e\,(137)(42)
    \left(1-\frac{D}{137N}\right){\rm eV}
    =
    0.07523660444808945~{\rm MeV}.
\end{equation}
Геометрический множитель:
\begin{equation}
    1+\frac{42}{137D}\left(1-\frac{13}{117}\right)
    =
    1.010481003181733.
\end{equation}

\begin{table}[h]
\centering
\caption{Аудит дефектов масс.  Residual считается по массе ядра.}
\begin{tabular}{lrrrr}
\toprule
Изотоп & $Z$ & Табличный дефект MeV & Модельный дефект MeV & Residual ppm \\
\midrule
$^1{\rm H}$     & 1  & $0.0000133172$ & $0.0000137339$ & $-0.000444$ \\
$^4{\rm He}$    & 2  & $28.2956897122$ & $0.3041006382$ & $7509.723753$ \\
$^{12}{\rm C}$  & 6  & $92.1618167441$ & $2.7369057434$ & $8002.327108$ \\
$^{56}{\rm Fe}$ & 26 & $492.259569845$ & $51.3930078482$ & $8463.590834$ \\
\bottomrule
\end{tabular}
\end{table}

Для железа:
\begin{equation}
    Z_{\rm Fe}=26=D,
\end{equation}
но совпадение номера не дает точной массы:
\begin{equation}
    \Delta M_{\rm Fe}^{\rm model}=51.393007848156245~{\rm MeV},
\end{equation}
\begin{equation}
    \Delta M_{\rm Fe}^{\rm ref}=492.2595698447403~{\rm MeV},
\end{equation}
\begin{equation}
    \epsilon_{\rm Fe}=8463.590833506674~{\rm ppm}.
\end{equation}
Критерий $<10$ ppm не выполнен.  Этот уровень фиксируется как отрицательный
результат.

\subsection{Уровень 2: eV, КЭД и лептонные массы}

Якорь:
\begin{equation}
    m_e=0.51099895069~{\rm MeV},
    \qquad
    \alpha^{-1}=137.035999084.
\end{equation}
Базовый шаг:
\begin{equation}
    5e~{\rm eV}=13.591409142295225~{\rm eV}.
\end{equation}
Для мюона:
\begin{equation}
    \frac{m_\mu}{m_e}
    =
    \alpha^{-1}\left(\frac{3}{2}+q\right)
    =
    206.768221426689.
\end{equation}
Для тау:
\begin{equation}
    \frac{m_\tau}{m_e}
    =
    \alpha^{-1}\left(N^2+qf_\tau\right),
    \qquad
    f_\tau=42+2e^{-2}+2e^{-7}.
\end{equation}
Получается
\begin{equation}
    \frac{m_\tau}{m_e}=3477.2282035579738.
\end{equation}
Четвертый заряженный массовый ориентир:
\begin{equation}
    \frac{m_{\tau'}}{m_e}
    =
    \alpha^{-1}\left[D^2+q(DN)\right]
    =
    92794.18434487357,
\end{equation}
\begin{equation}
    m_{\tau'}=47.417730830364825~{\rm GeV}.
\end{equation}
Это не разрешенный обычный последовательный четвертый лептон.

Стерильный нейтральный массовый ориентир:
\begin{equation}
    m_{\nu_{\tau'}}
    =
    m_e
    \left(\frac{F_{26}}{\alpha^{-1}}\right)
    \left[
        1+\frac{DN}{\pi}\delta\phi(1-s)
    \right],
\end{equation}
\begin{equation}
    m_{\nu_{\tau'}}=1.5566463427697075~{\rm GeV}.
\end{equation}
Это только масса.  Реликтовая плотность, время жизни и смешивание требуют
отдельной модели.

\subsection{Уровень 3: молекулярный уровень, восстановление и семантика}

Электронная структура молекул здесь не считается.  На этом уровне есть только
символьный канал восстановления:
\begin{equation}
    A=\mathbf{1}_{R\ge42}.
\end{equation}
Если повреждено не более $10$ осей, остается большинство:
\begin{equation}
    d_{\rm damaged}\le10
    \quad\Rightarrow\quad
    26-d_{\rm damaged}\ge16.
\end{equation}
Тестовая матрица проверяет это для синтетической ДНК и семантического архива:
\begin{equation}
    120{,}000~{\rm tokens}
    \mapsto
    117~{\rm vectors}
    \mapsto
    117\times130~{\rm bytes}.
\end{equation}
Это корректирующий код в тестовой модели, а не биохимическая теория.

\subsection{Уровень 4: макроуровень, газ и термодинамический предел}

Пусть микросостояние хранит вычет $a\in\mathbb{F}_{137}$, который отображается
в центрированный целый диапазон
\begin{equation}
    \rho(a)\in\{-68,-67,\ldots,67,68\}.
\end{equation}
Для трех компонент скорости:
\begin{equation}
    v=(\rho(a_1),\rho(a_2),\rho(a_3))v_0,
\end{equation}
\begin{equation}
    \varepsilon(v)=\frac{m}{2}|v|^2.
\end{equation}
Максимизация энтропии
\begin{equation}
    S=-k_B\sum_i p_i\ln p_i
\end{equation}
при условиях
\begin{equation}
    \sum_i p_i=1,\qquad \sum_i p_i\varepsilon_i=\bar{E}
\end{equation}
дает
\begin{equation}
    p_i=\frac{e^{-\beta\varepsilon_i}}{Z_1(\beta)}.
\end{equation}
В изотропном пределе сумма по вычетам заменяется интегралом по скоростям:
\begin{equation}
    Z_{N_g}
    =
    \frac{1}{N_g!}
    \left[
        V\left(\frac{2\pi m}{\beta h^2}\right)^{3/2}
    \right]^{N_g}.
\end{equation}
Давление:
\begin{equation}
    P=\frac{1}{\beta}\frac{\partial \ln Z_{N_g}}{\partial V}
      =
    \frac{N_g}{\beta V}.
\end{equation}
Так как $\beta^{-1}=k_BT$,
\begin{equation}
    PV=N_gk_BT=nRT.
\end{equation}
Энтропия:
\begin{equation}
    S=k_B\ln\Omega.
\end{equation}
GF(137) здесь является конечной регуляризацией микросостояний.  Закон
идеального газа получается в стандартном пределе усреднения.

\subsection{Уровень 5: космологический масштаб}

Референсные значения:
\begin{equation}
    H_{0,\rm Planck}=67.4,
    \qquad
    H_{0,\rm SH0ES}=73.04.
\end{equation}
Доля расхождения:
\begin{equation}
    \frac{73.04-67.4}{67.4}=0.0836795252225519.
\end{equation}
CPL-пара:
\begin{equation}
    w_0=-0.7519948527423932,
    \qquad
    w_a=-0.8600649695978589.
\end{equation}

Проверенный шлюз времени:
\begin{equation}
    H_{\rm gate}(z,x)
    =
    \left[
        1-\frac{q}{117}e^z\cos^2(5x)
    \right]^{-1/2}.
\end{equation}
Численно:
\begin{equation}
    \frac{q}{117}=7.57317\times10^{-5},
    \qquad
    \max H_{\rm gate}\simeq1.00028.
\end{equation}
Лучший принятый ряд:
\begin{equation}
    w_0=-0.920331,\qquad w(z=1)=-1.00259.
\end{equation}
Цель:
\begin{equation}
    w_{\rm target}(z=1)=-1.18202733754.
\end{equation}
Остаток:
\begin{equation}
    w(z=1)-w_{\rm target}(z=1)=0.179442.
\end{equation}
Это отрицательный результат на DESI-масштабах.

CMB-инвариант:
\begin{equation}
    \Delta N_{\rm eff}
    =
    \frac{D}{137}\delta\phi(1-q)
    =
    0.011098917626279356.
\end{equation}
\begin{equation}
    N_{\rm eff}=3.0570989176262793.
\end{equation}
Относительно Planck 2018 CMB+BAO $2.99\pm0.17$ это
\begin{equation}
    0.39469951544870036
\end{equation}
стандартных отклонений.

Для массы
\begin{equation}
    m_{\nu_{\tau'}}=1.5566463427697075~{\rm GeV}
\end{equation}
оценка cutoff дает
\begin{equation}
    \lambda_{\rm cut}
    =
    6.097325934119188\times10^{-7}~{\rm kpc}.
\end{equation}
Это ниже $100~{\rm kpc}$ и ниже DESI/BAO-масштабов.  Вывод условен: нужна
стерильная, холодная или достаточно нетепловая история рождения.

\section{Прикладная целочисленная реализация}

GF(137)-инференс хранит веса как байты:
\begin{equation}
    M_{\rm GF(137)}=1.62793~{\rm KB},
    \qquad
    M_{\rm float32}=6.50391~{\rm KB}.
\end{equation}
Это примерно $4\times$ по памяти.  На $1000$ прямых проходов:
\begin{equation}
    \frac{58.7784}{12.1055}\simeq4.86.
\end{equation}
C++ путь через внешний вызов дает $40.1530$ ms и ускорение около $1.46\times$.

\begin{table}[h]
\centering
\caption{Аудит памяти и времени для edge inference.}
\small
\setlength{\tabcolsep}{3pt}
\begin{tabularx}{\linewidth}{@{}YrrP{0.18\linewidth}@{}}
\toprule
Реализация & Память KB & Время для $1000$ проходов & Роль \\
\midrule
float32 NumPy MLP & 6.50391 & 58.7784 ms & базовая линия \\
GF(137) C++ fused \texttt{uint8\_t} & 1.62793 & 40.1530 ms & внешний вызов \\
GF(137) C++ NEON native loop & 1.62793 & 12.1055 ms & нативный цикл \\
\bottomrule
\end{tabularx}
\end{table}

Аппаратные примитивы:
\begin{align}
    {\tt mac137}(a,b,c)&=(c+ab)\bmod137,\\
    {\tt inv137}(a)&=a^{-1}\bmod137,\quad a\ne0,\\
    {\tt ge42}(a)&=\mathbf{1}_{a\ge42}.
\end{align}
Преимущество здесь не в новой математике, а в ограниченной целочисленной
арифметике без общего делителя в горячем цикле.

\section{Сводная таблица}

\begin{table}[h]
\centering
\caption{Пятиуровневый статус.}
\small
\setlength{\tabcolsep}{3pt}
\begin{tabularx}{\linewidth}{@{}P{0.10\linewidth}P{0.18\linewidth}P{0.34\linewidth}Y@{}}
\toprule
Уровень & Масштаб & Основная величина & Статус \\
\midrule
1 & MeV/КХД & остаток Fe $8463.590833506674$ ppm & не проходит $<10$ ppm \\
2 & eV/КЭД & $m_{\nu_{\tau'}}=1.5566463427697075$ GeV & массовый ориентир \\
3 & мол./сем. & $26$ осей, $10$ повреждений & тест восстановления проходит \\
4 & газ/тепловой & $PV=N_gk_BT=nRT$ & предел усреднения \\
5 & космология & $N_{\rm eff}=3.0570989176262793$ & в окне Planck \\
\bottomrule
\end{tabularx}
\end{table}

\section{Ограничения}

Отрицательный ядерный результат сохраняется как часть вывода.  В текст не
добавляется подгоночный член, который принудительно согласовал бы NIST/CODATA
массы.  Лептонные формулы являются параметризациями, а не выводом из
лагранжиана.  Молекулярный и DNA-разделы описывают символьный код
восстановления, а не биохимию.  Газовый раздел -- стандартная статистическая
механика с конечным алфавитом микросостояний.  Космологический раздел
фиксирует проверки замыкания и отрицательные результаты при явно указанных
предположениях.

\section{Заключение}

Конструкция e-5-137 полезна в этом репозитории как компактный язык
аудита.  Она дает воспроизводимые целочисленные ядра, фиксированные пороги
восстановления, несколько численных массовых ориентиров и явные отрицательные
результаты.  Пятиуровневая структура отделяет инженерные тесты от физических
утверждений, которые пока не подтверждены.  Основной критерий -- таблица
остатков, а не совпадение целых чисел.

\begin{thebibliography}{11}

\bibitem{pdg2024}
S. Navas et al. (Particle Data Group),
``Review of Particle Physics,''
\emph{Phys. Rev. D} \textbf{110}, 030001 (2024), with 2025 update.

\bibitem{adhikari2017}
R. Adhikari et al.,
``A White Paper on keV Sterile Neutrino Dark Matter,''
\emph{JCAP} \textbf{01}, 025 (2017).

\bibitem{planck2018}
N. Aghanim et al. (Planck Collaboration),
``Planck 2018 results. VI. Cosmological parameters,''
\emph{Astron. Astrophys.} \textbf{641}, A6 (2020),
arXiv:1807.06209.

\bibitem{actdr6extended}
Atacama Cosmology Telescope Collaboration,
``The Atacama Cosmology Telescope: DR6 Constraints on Extended Cosmological
Models,''
arXiv:2503.14454.

\bibitem{nistatomicweights}
National Institute of Standards and Technology,
``Atomic Weights and Isotopic Compositions with Relative Atomic Masses,''
\url{https://www.nist.gov/pml/atomic-weights-and-isotopic-compositions}.

\bibitem{codata2018}
National Institute of Standards and Technology,
``CODATA Recommended Values of the Fundamental Physical Constants: 2018,''
\url{https://pml.nist.gov/cuu/pdf/wall_2018.pdf}.

\bibitem{kuratov2025}
Y. Kuratov, M. Arkhipov, A. Bulatov, and M. Burtsev,
``Cramming 1568 Tokens into a Single Vector and Back Again: Exploring the
Limits of Embedding Space Capacity,''
arXiv:2502.13063 (2025).

\bibitem{boltzmann1896}
L. Boltzmann,
\emph{Vorlesungen \"uber Gastheorie},
J. A. Barth, Leipzig (1896).

\bibitem{gibbs1902}
J. W. Gibbs,
\emph{Elementary Principles in Statistical Mechanics},
Charles Scribner's Sons, New York (1902).

\end{thebibliography}

\end{document}
